\chapter{应用层}

\section{应用层协议}

\subsection{网络应用}

网络应用程序是指能够运行在不同端系统并通过网络彼此通信的程序。例如在Web应用程序中，有两个互相通信的程序，一个是运行在用户主机上的浏览器程序，另一个是运行在Web服务器主机上的服务器程序。\\

应用程序体系结构包括Client-Server和P2P（Peer-to-Peer）两种体系结构。在Client-Server中，服务器主机一直运行处理客户端的请求。因为服务器需要有一个固定的的IP地址，因此客户端能够通过IP地址发送分组与其联系。著名的Client-Server应用程序包括Web、FTP、Telnet和电子邮件等。在Client-Server应用中，经常会出现一台单独的服务器主机不堪重负，跟不上所有客户请求的情况。例如搜索引擎（如Google、Bing、Baidu）、电商（如Amazon、eBay、Alibaba）、电子邮件（如Gmail、Yahoo!）、社交媒体（如Facebook、Instagram、Twitter、WeChat）都分步了多个数据中心，每数据中心都有数十万台服务器，必须要持续供电和维护。而在P2P中，主机之间的通信不必通过专门的服务器，很多流量密集型的应用都是P2P体系结构的，包括对等方协助下载加速器（如迅雷）、文件共享（如BitTorrent）。\\

\subsection{进程通信}

在同一个端系统中，多个进程可以通过进程间通信的机制相互通信，然而在不同的端系统（可能有不同的操作系统）中，需要通过网络交换报文（message）。进程通过套接字（socket）接口向网络发送或接收报文。\\

\begin{figure}[H]
    \centering
    \begin{tikzpicture}
        \draw (0,0) rectangle (3,9);
        \draw (8,0) rectangle (11,9);
        \draw (0,1) -- (3,1);
        \draw (0,2) -- (3,2);
        \draw (0,3) -- (3,3);
        \draw (0,4) -- (3,4);
        \draw (0,5) -- (3,5);
        \draw (0,6.5) -- (3,6.5);
        \draw (8,1) -- (11,1);
        \draw (8,2) -- (11,2);
        \draw (8,3) -- (11,3);
        \draw (8,4) -- (11,4);
        \draw (8,5) -- (11,5);
        \draw (8,6.5) -- (11,6.5);

        \draw (1.5,8.5) node{应用层};
        \draw (1.5,5.5) node{表示层};
        \draw (1.5,4.5) node{会话层};
        \draw (1.5,3.5) node{传输层};
        \draw (1.5,2.5) node{网络层};
        \draw (1.5,1.5) node{数据链路层};
        \draw (1.5,0.5) node{物理层};
        \draw (9.5,8.5) node{应用层};
        \draw (9.5,5.5) node{表示层};
        \draw (9.5,4.5) node{会话层};
        \draw (9.5,3.5) node{传输层};
        \draw (9.5,2.5) node{网络层};
        \draw (9.5,1.5) node{数据链路层};
        \draw (9.5,0.5) node{物理层};

        \draw (1.5,7.5) ellipse (1 and 0.5);
        \draw (9.5,7.5) ellipse (1 and 0.5);
        \draw (1.5,7.5) node{进程};
        \draw (9.5,7.5) node{进程};

        \draw[orange, very thick] (0.75,6) rectangle (2.25,6.75);
        \draw[orange, very thick] (1.5,6.25) node{socket};
        \draw[orange, very thick] (8.75,6) rectangle (10.25,6.75);
        \draw[orange, very thick] (9.5,6.25) node{socket};

        \draw[->, very thick] (-1,9) -- (-1,0);
        \draw[->, very thick] (12,0) -- (12,9);
        \draw[->] (1.5,0) -- (1.5,-1) -- (9.5,-1) -- (9.5,0);
    \end{tikzpicture}
    \caption{socket}
\end{figure}

\vspace{0.5cm}

在进程发送分组的过程中，必须要标识IP地址和端口号（port number）才能将分组发送给另一主机的进程。其中IP地址用于标识主机，端口号用于指定运行在目的主机上的接收进程。由于一台主机上会运行多个应用程序，因此端口号是不可或缺的信息。一些著名的应用已经被分配了特定的端口号，例如Web服务器使用端口号80、邮件服务器（SMTP协议）进程使用端口号25。\\

发送端在使用socket时必须选择一种传输层协议，不同的协议会提供不同的服务。\\

一个传输层协议可以通过四个方面进行分类：

\begin{enumerate}
    \item 可靠数据传输（reliable data transfer）：分组在网络传输中可能会因溢出或损坏等原因丢失，对于电子邮件、文件传输和金融相关的应用来说，数据丢失会造成灾难性的后果，这种情况下就必须采用可靠数据传输。对于一些可以容忍丢失（loss-tolerant）的应用，例如多媒体音视频，它们能够承受一定量的数据丢失，这只会造成小干扰，而非致命性的问题。

    \item 吞吐量（throughput）：传输层协议能够以特定的速率提供服务。

    \item 定时（timing）：传输层协议提供定时保证。例如在网络电话或多人游戏中，较长的时间延迟会出现不自然的停顿或失去真实感。

    \item 安全性（security）：传输层协议保证数据安全。例如在发送端将数据加密，并在接收端解密数据，以防在传输被中途被窃听。
\end{enumerate}

\begin{table}[H]
    \centering
    \setlength{\tabcolsep}{5mm}{
        \begin{tabular}{|c|c|c|c|}
            \hline
            \textbf{应用} & \textbf{数据丢失} & \textbf{吞吐量}           & \textbf{时间敏感} \\
            \hline
            文件传输      & 不允许丢失        & 弹性                      & 不                \\
            \hline
            电子邮件      & 不允许丢失        & 弹性                      & 不                \\
            \hline
            网络电话      & 容忍丢失          & few kbps $ \sim $ 1 Mbps  & 100 ms            \\
            \hline
            视频会议      & 容忍丢失          & 10 kbps $ \sim $ 5 Mbps   & 100 ms            \\
            \hline
            交互式游戏    & 容忍丢失          & few kbps $ \sim $ 10 kbps & 100 ms            \\
            \hline
        \end{tabular}
    }
    \caption{常见应用传输服务需求}
\end{table}

\vspace{0.5cm}

\subsection{TCP / UDP}

TCP （Transmission Control Protocol）的特点包括：

\begin{itemize}
    \item 面向连接服务（connection-oriented service）：在应用层数据包开始发送之前，TCP让客户端和服务器之间互相交换传输层控制信息，让它们为分组的到来做好准备。在此之后，两个进程的socket就能建立起TCP连接，并可以发送报文了。在发送结束后，该TCP连接会被拆除。

    \item 可靠数据传输：TCP确保了通信进程交付的数据无差错、不丢失、不重复、不乱序。

    \item 拥塞控制（congestion control）
\end{itemize}

UDP（User Datagram Protocol）的特点包括：

\begin{itemize}
    \item 无连接（connectionless）

    \item 不可靠数据传输：UDP不保证报文能够到达接收端，同时报文也有可能是乱序到达的。

    \item 无拥塞控制
\end{itemize}

\begin{table}[H]
    \centering
    \setlength{\tabcolsep}{5mm}{
        \begin{tabular}{|c|c|}
            \hline
            \textbf{应用} & \textbf{传输协议} \\
            \hline
            文件传输      & TCP               \\
            \hline
            电子邮件      & TCP               \\
            \hline
            网络电话      & UDP               \\
            \hline
            视频会议      & UDP               \\
            \hline
            交互式游戏    & UDP               \\
            \hline
        \end{tabular}
    }
    \caption{常见应用传输协议}
\end{table}

\newpage

\section{HTTP}

\subsection{HTTP （HyperText Transfer Protocol）}

一个Web页面是由对象（object）组成的，一个对象就是一个文件，例如HTML文件、JPEG图片、JavaScript文件、CSS样式文件等。如果一个Web页面包含1个HTML文件和5个JPEG图片，那么这个Web页面就有6个对象。\\

每一个对象都可以通过URL（Uniform Resource Locator）寻址，URL地址由存放对象的服务器主机名（host name）和路径名（path name）组成。例如对于http://www.someSchool.edu/someDepartment/picture.gif而言，其中主机名就是www.someSchool.edu，路径名是/someDepartment/picture.gif。\\

HTTP定义在RFC 1945、RFC 7230和RFC 7540中。HTTP使用TCP作为它的支撑传输协议，客户端首先发起TCP连接，连接建立后，客户进程可以向服务器进程发送HTTP请求报文，服务器进程可以向客户进程发送HTTP响应报文。\\

HTTP是一个无状态协议（stateless protocol），服务器不能存储任何关于客户的状态信息。例如某个客户在短时间内两次请求同一个对象，服务器并不会因为第一次已经向客户提供了该对象而不再作出响应，而是再次重新发送对象。\\

\subsection{HTTP请求报文}

HTTP请求报文由第一行的请求行（request line）和后续的首部行（header line）组成，每行由一个回车（carriage return）和换行（line feed）结束，在首部行之后再附加一个只包含回车换行的空行 。\\

\mybox{HTTP请求报文}

\begin{lstlisting}
GET /somedir/page.html HTTP/1.1\r\n
Host: www.someschool.edu\r\n
Connection: close\r\n
User-agent: Mozilla/5.0\r\n
Accept-language: fr\r\n
\r\n
[entity body]
\end{lstlisting}

请求行包含3个字段：

\begin{itemize}
    \item 方法：包括GET、POST、HEAD、PUT和DELETE。

    \item URL：带有请求对象的标识。

    \item HTTP版本
\end{itemize}

首部行中Host指明了对象所在的主机。Connection: close表示浏览器告诉服务器不要使用持续连接，而是要求服务器在发送完对象后就关闭此连接。User-agent指明了用户代理，即向服务器发送请求的浏览器类型，例如Mozilla/5.0。服务器可以通过此信息向用户代理发送不同版本的对象。Accept-language表示用户想要获取对象的语言版本，如果服务器不存在的话，就会发送一个默认版本。\\

使用GET方法时，实体部分（entity body）为空。而使用POST方法时，实体部分可以用于包含用户在表单中填写的输入值。HEAD方法与GET类似，服务器会响应，但并不返回请求对象。因此HEAD方法常用于调试跟踪。PUT方法允许用户向服务器上传对象。DELETE方法允许用户删除服务器上的对象。\\

\subsection{HTTP响应报文}

HTTP响应报文由状态行（status line）、首部行（header line）和实体部分组成。\\

\mybox{HTTP响应报文}

\begin{lstlisting}
HTTP/1.1 200 OK\r\n
Connection: close\r\n
Date: Tue, 18 Aug 2015 15:44:04 GMT\r\n
Server: Apache/2.2.3 (CentOS)\r\n
Last-Modified: Tue, 18 Aug 2015 15:11:03 GMT\r\n
Content-Length: 6821\r\n
Content-Type: text/html\r\n
\r\n
(data data data data data ...)
\end{lstlisting}

状态行包括了协议版本、状态码和对应状态信息。首部行中Connection: close用于告诉客户，发送完报文后将会关闭TCP连接。Date用于表示服务器生成并发送该报文的日期时间。Last-Modified用于表示对象创建或最后修改的日期时间。Content-Length用于表示被发送对象的字节数。Content-Type用于表示对象的类型。\\

\begin{table}[H]
    \centering
    \setlength{\tabcolsep}{5mm}{
        \begin{tabular}{|l|l|}
            \hline
            \textbf{状态码}                & \textbf{含义}                        \\
            \hline
            200 OK                         & 请求成功                             \\
            \hline
            204 No Content                 & 无内容                               \\
            \hline
            301 Moved Permanently          & 永久性重定向，资源被分配了新URL      \\
            \hline
            400 Bad Request                & 请求语法错误，服务器无法理解         \\
            \hline
            403 Forbidden                  & 拒绝执行请求                         \\
            \hline
            404 Not Found                  & 无法找到资源                         \\
            \hline
            500 Internal Server Error      & 服务器内部错误                       \\
            \hline
            503 Service Unavailable        & 由于超载或系统维护，暂时无法处理请求 \\
            \hline
            505 HTTP Version not supported & 不支持请求的HTTP协议的版本           \\
            \hline
        \end{tabular}
    }
    \caption{常见HTTP状态码}
\end{table}

\newpage

\section{Cookie}

\subsection{Cookie}

HTTP服务器是无状态的，然而一些Web站点通常希望能够识别用户，从而记住用户信息或限制用户访问。目前Cookie广泛用于记录用户登录信息，这样下次访问时可以不再输入用户名和密码了。当然这种方便也存在用户信息泄密的问题，尤其在多个用户公用一台电脑时很容易出现这样的情况。\\

\begin{figure}[H]
    \centering
    \includegraphics[scale=0.15]{img/C2/2-3/1.png}
\end{figure}

例如用户首次访问Amazon，HTTP响应报文中会包含Set-cookies识别码，浏览器会将识别码添加到所管理的文件中。当用户继续浏览Amazon时，每一个HTTP请求浏览器都会从cookie文件中查询该网站的识别码，并添加到HTTP请求报文中。Amazon也可以通过cookie来维护用户希望购买的商品信息，并推荐个性化产品。

\begin{figure}[H]
    \centering
    \begin{tikzpicture}
        \node at (0,15) {Client};
        \node at (8,15) {Server};

        \draw[loosely dashed, <-] (0,0) -- (0,14);
        \draw[loosely dashed, <-] (8,0) -- (8,14);

        \draw[blue, ->, thick] (0,14) -- (8,12) node[above, midway, sloped]{HTTP request};
        \draw[blue, ->, thick] (8,12) -- (0,10) node[above, midway, sloped]{HTTP response (Set-cookie: 1678)};
        \draw[blue, ->, thick] (0,10) -- (8,8) node[above, midway, sloped]{HTTP request (cookie: 1678)};
        \draw[blue, ->, thick] (8,8) -- (0,6) node[above, midway, sloped]{HTTP response};
        \draw[blue, ->, thick] (0,5) -- (8,3) node[above, midway, sloped]{HTTP request (cookie: 1678)};
        \draw[blue, ->, thick] (8,3) -- (0,1) node[above, midway, sloped]{HTTP response};

        \node at (-2,14) [cylinder, shape border rotate=90, draw, minimum height=12mm, minimum width=10mm] {};
        \node at (-2,13) {eBay: 8734};

        \node at (-2,10) [cylinder, shape border rotate=90, draw, minimum height=12mm, minimum width=10mm] {};
        \node at (-2,9) {eBay: 8734};
        \node at (-2,8.5) {Amazon: 1678};

        \node at (-2,5) [cylinder, shape border rotate=90, draw, minimum height=12mm, minimum width=10mm] {};
        \node at (-2,4) {eBay: 8734};
        \node at (-2,3.5) {Amazon: 1678};

        \node at (11,7) [cylinder, shape border rotate=90, draw, minimum height=12mm, minimum width=10mm] {};
        \node at (11,6) {database};

        \draw[dotted, ->, thick] (8,12) -- (11,7.7) node[above, midway, sloped]{creates ID 1678};
        \draw[dotted, ->, thick] (8,8) -- (10.5,7.5) node[above, midway, sloped]{access};
        \draw[dotted, ->, thick] (8,3) -- (10.5,7) node[above, midway, sloped]{access};
    \end{tikzpicture}
    \caption{Cookie}
\end{figure}

\newpage

\section{Email}

\subsection{SMTP（Simple Mail Transfer Protocol）}

电子邮件系统由用户代理（user agent）、邮件服务器（mail server）和SMTP组成，其中例如Outlook、Apple Mail、Gmail就是用户代理。当Alice发送邮件时，她的用户代理就会向邮件服务器发送邮件，此时邮件被放在邮件服务器的待发报文队列中。当Bob要查看邮件时，他的用户代理就会其邮件服务器中获取该报文。\\

\begin{figure}[H]
    \centering
    \begin{tikzpicture}
        \draw (0,0) rectangle (2.5,2);
        \node at (1.25,1.25) {Alice's};
        \node at (1.25,0.75) {User Agent};

        \draw (4,0) rectangle (6.5,2.5);
        \node at (5.25,2) {Alice's};
        \node at (5.25,1.5) {Mail Server};
        \draw (4.5,0.75) rectangle (6,1);
        \draw (4.3,0.25) rectangle (4.5,0.5);
        \draw (4.7,0.25) rectangle (4.9,0.5);
        \draw (5.1,0.25) rectangle (5.3,0.5);
        \draw (5.5,0.25) rectangle (5.7,0.5);
        \draw (5.9,0.25) rectangle (6.1,0.5);

        \draw (8,0) rectangle (10.5,2.5);
        \node at (9.25,2) {Bob's};
        \node at (9.25,1.5) {Mail Server};
        \draw (8.5,0.75) rectangle (10,1);
        \draw (8.3,0.25) rectangle (8.5,0.5);
        \draw (8.7,0.25) rectangle (8.9,0.5);
        \draw (9.1,0.25) rectangle (9.3,0.5);
        \draw (9.5,0.25) rectangle (9.7,0.5);
        \draw (9.9,0.25) rectangle (10.1,0.5);

        \draw (12,0) rectangle (14.5,2);
        \node at (13.25,1.25) {Bob's};
        \node at (13.25,0.75) {User Agent};

        \draw[->, thick] (2.5,1) -- (4.5,1);
        \draw[->, thick] (6,1) -- (8.25,0.5);
        \draw[->, thick] (10,0.5) -- (12,1);
    \end{tikzpicture}
    \caption{SMTP}
\end{figure}

SMTP是电子邮件中使用的应用层协议，它使用TCP可靠数据传输服务，使用25号端口。\\

SMTP使用命令和应答对报文进行传输：\\

\begin{table}[H]
    \centering
    \setlength{\tabcolsep}{5mm}{
        \begin{tabular}{|c|l|}
            \hline
            \textbf{命令} & \textbf{描述}        \\
            \hline
            HELO          & 与服务器确认信息     \\
            \hline
            AUTH          & 登录验证             \\
            \hline
            MAIL FROM     & 发件人信息           \\
            \hline
            RCPT TO       & 收信人地址           \\
            \hline
            DATA          & 输入邮件基本信息     \\
            \hline
            FROM          & 发信人显示信息       \\
            \hline
            TO            & 服务器收件人显示信息 \\
            \hline
            SUBJECT       & 邮件标题             \\
            \hline
            QUIT          & 断开链接             \\
            \hline
        \end{tabular}
    }
    \caption{SMTP命令}
\end{table}

\begin{table}[H]
    \centering
    \setlength{\tabcolsep}{5mm}{
        \begin{tabular}{|c|l|}
            \hline
            \textbf{返回码} & \textbf{含义}    \\
            \hline
            220             & 服务就绪         \\
            \hline
            250             & 请求动作成功完成 \\
            \hline
            235             & 认证通过         \\
            \hline
            221             & 处理中           \\
            \hline
            354             & 发送开始         \\
            \hline
            500             & 指令错误         \\
            \hline
            550             & 命令无法执行     \\
            \hline
        \end{tabular}
    }
    \caption{SMTP应答返回码}
\end{table}

\vspace{0.5cm}

当用户调用用户代理查看邮件报文时就需要用到邮件访问协议。流行的邮件访问协议有邮局协议版本3（POP3）、因特网邮件访问协议（IMAP, Internet Mail Access Protocol）和HTTP。

\newpage